% Lec 01–02：命题逻辑与一阶逻辑入门
\documentclass[UTF8,a4paper,11pt]{ctexart}
\usepackage[a4paper,margin=2.3cm]{geometry}
\usepackage{amsmath,amssymb,mathtools}
\usepackage{enumitem}
\usepackage[hidelinks]{hyperref}
\usepackage{fancyhdr}
% ctexart does not provide semantic double brackets by default.
\newcommand{\llbracket}{\left[\!\left[}
\newcommand{\rrbracket}{\right]\!\right]}
\setlist[itemize]{leftmargin=2.2em,itemsep=0.25em,topsep=0.35em}
\setlength{\parindent}{2em}
\setlength{\parskip}{0.35em}
\pagestyle{fancy}\fancyhf{}
\fancyhead[L]{离散数学与结构（I）}\fancyhead[R]{Lec 01–02}
\fancyfoot[C]{\thepage}\setlength{\headheight}{14pt}
\title{命题逻辑与一阶逻辑入门}
\author{zhangjunshuo}\date{2026 年 9 月}
\begin{document}
\maketitle
\begin{quote}
本人选课较晚，Lec 01–02 主要依据作业与参考资料整理，非课堂记录。本稿整理自然演绎、Fitch 列表、子证明作用域、经典相继式和一阶逻辑基础。
\end{quote}
\tableofcontents\clearpage
\section{自然演绎与逆否式}


\begin{quote}
来源说明：本人选课较晚，本节主要依据作业与参考教材整理，非课堂记录。
\end{quote}


\subsection{记号与两种“反证”动作}


采用课程约定：$\neg\varphi$ 是 $\varphi\to\bot$ 的缩写。

\begin{itemize}
\item \textbf{$\to I$}：在临时假设 $[\varphi]^i$ 下推出 $\psi$，解除该假设得到 $\varphi\to\psi$。若推出的是 $\bot$，得到 $\neg\varphi$。
\item \textbf{经典 RAA}：在临时假设 $[\neg\varphi]^i$ 下推出 $\bot$，解除它得到 $\varphi$。
\item \textbf{$\to E$}：由 $\varphi\to\psi$ 和 $\varphi$ 得 $\psi$。
\end{itemize}


所以“假设 $\varphi$ 导出矛盾”得到的是 $\neg\varphi$，而“假设 $\neg\varphi$ 导出矛盾”才使用 RAA 得到 $\varphi$。

\subsection{逆否式}


证明

\[
(\neg Q\to\neg P)\leftrightarrow(P\to Q).
\]

正向推导：

\begin{verbatim}
1.  ¬Q→¬P        AS
2.    P           AS
3.      ¬Q        AS
4.      ¬P        →E 1,3
5.      ⊥         →E 2,4
6.    Q           RAA 3
7.  P→Q           →I 2-6
8.  (¬Q→¬P)→(P→Q) →I 1-7
\end{verbatim}


反向推导：

\begin{verbatim}
1.  P→Q           AS
2.    P           AS
3.    Q           →E 1,2
4.      ¬Q        AS
5.      ⊥         →E 3,4
6.    ¬P          →I 2-5
7.  ¬Q→¬P         →I 4-6
8.  (P→Q)→(¬Q→¬P) →I 1-7
\end{verbatim}


编号只表示要解除的假设。第一段的 RAA 解除第 3 行的 $[\neg Q]$；第二段的第 6 行是由 $[P]$ 导出矛盾，所以必须写 $\to I$。

\subsection{Fitch 列表与 Reiteration}


典型证明：

\begin{verbatim}
1 | ¬¬P        PR
2 |   ¬P       AS
3 |   ⊥        →E 2,1
4 | P          RAA 2
\end{verbatim}


缩进表示临时假设的作用域。已经写过且仍在作用域中的公式可以用 \texttt{R, 行号} 重复引用；已关闭子证明中的行不能再引用。
\clearpage
\section{子证明作用域与经典相继式}


\subsection{作用域}


条件消去是

\begin{verbatim}
m | A→B
n | A
  | B          →E m,n
\end{verbatim}


可以开启额外假设：

\begin{verbatim}
1 | R              PR
2 |   L            AS
3 |   R∧L          ∧I 1,2
4 | L→(R∧L)        →I 2-3
\end{verbatim}


第 4 行不再依赖临时假设 $L$，但仍依赖第 1 行的 $R$。关闭子证明后，任何依赖其假设的行都不能在外层继续使用；这正是防止从矛盾假设任意推出结论的作用域条件。

\subsection{相继式的语义}


相继式写作 $\Gamma\Rightarrow\Delta$。这里 $\Gamma,\Delta$ 按有限多重集理解，结构规则通常默认可用。它在赋值 $v$ 下为真，当且仅当

\[
\bigl(v\models\gamma\text{ 对所有 }\gamma\in\Gamma\bigr)
\Longrightarrow
\bigl(v\models\delta\text{ 对某个 }\delta\in\Delta\bigr).
\]

\subsection{八类联结词规则}


下面只写主公式，未写出的上下文在前后两行保持不变：

\begin{itemize}
\item $\neg\Rightarrow$：把左侧的 $\neg\varphi$ 移到右侧的 $\varphi$；$\Rightarrow\neg$：把右侧的 $\neg\varphi$ 移到左侧的 $\varphi$。
\item $\land\Rightarrow$：左侧的 $\varphi\land\psi$ 拆成两个左侧公式；$\Rightarrow\land$：右侧合取需要分别证明两个支。
\item $\Rightarrow\lor$：右侧析取可展开为两个右侧公式；$\lor\Rightarrow$：左侧析取需要两个分支。
\item $\Rightarrow\to$：证明右侧 $\varphi\to\psi$ 时，把 $\varphi$ 放到左侧、$\psi$ 放到右侧；$\to\Rightarrow$：左侧蕴含需要两个分支。
\end{itemize}


公理模式是 $\Gamma\Rightarrow\Delta$ 且 $\Gamma\cap\Delta\ne\varnothing$。无 cut 推导满足子公式性；做题时通常从目标相继式自下而上选择主公式。

本节的 System K 是手写笔记中的\textbf{经典相继式系统}。25–26 讲义后面另有模态逻辑 System K，并使用 $\Box,\Diamond$；不能把两者的规则混在一起。
\clearpage
\section{一阶逻辑的语法、语义与量词}


\subsection{语法与代入}


项由变量、常元和函数符号归纳生成；公式由谓词原子式、等式原子式、联结词和量词归纳生成。没有自由变量的公式称为句子；没有变量的项称为闭项。

代入 $\varphi[t/x]$ 只替换 $x$ 的自由出现。若遇到 $Qx$，该辖域内的 $x$ 不再替换；若遇到 $Qy$ 且 $y\ne x$，还必须检查 $t$ 中的自由变量不会被新量词捕获。必要时先对约束变量作 $\alpha$-改名。

\subsection{结构与满足}


结构 $\mathcal M$ 给出非空论域 $\Omega$，并解释每个常元、函数符号和关系符号。变量赋值 $s$ 把变量送入论域；项值与公式真值按语法递归定义。

记 $\mathcal M,s\models\varphi$ 为带赋值满足，$\mathcal M\models\varphi$ 为句子满足，$\models\varphi$ 为所有结构中有效，$\Gamma\models\varphi$ 为语义蕴涵。等号按实际相等解释：

\[
\mathcal M,s\models t_1=t_2
\quad\Longleftrightarrow\quad
\llbracket t_1\rrbracket_{\mathcal M,s}=\llbracket t_2\rrbracket_{\mathcal M,s}.
\]

\subsection{量词规则}


存在量词可取为缩写 $\exists x\varphi:=\neg\forall x\neg\varphi$。核心规则是

\[
\frac{\Gamma\vdash\varphi}{\Gamma\vdash\forall x\varphi}\ (\forall I),
\qquad
\frac{\Gamma\vdash\forall x\varphi}{\Gamma\vdash\varphi[t/x]}\ (\forall E).
\]

\begin{itemize}
\item $\forall I$ 要求 $x$ 不在任何未解除假设中自由出现；
\item $\forall E$ 要求 $t$ 对 $x$ 可自由代入；
\item 在非空论域中，$\forall x\varphi(x)\vdash\exists x\varphi(x)$；
\item 从 $\varphi(t)$ 可推出 $\exists x\varphi(x)$。
\end{itemize}


例如 $P(x)\vdash\forall x\exists yP(y)$ 是合法的，但不能直接对开放假设中的 $x$ 使用 $\forall I$。合法路线是先得出与 $x$ 无关的 $\exists yP(y)$，再用新的变量 $z$ 做 $\forall I$，最后对绑定变量作 $\alpha$-改名。相反，从开放假设 $P(x)$ 直接得到 $\forall xP(x)$ 是无效的。

\end{document}
